\section{Background and Related Work}
\label{sec:background}
\label{sec:related}

% This is ONE combined Section II. Pick ONE name for it -- "Background
% and Related Work", "Background", or "Related Work" -- and stick with it.
%
% Do NOT split this into "Background" + "Related Work" (neither as two
% peer top-level sections, NOR as two subsections named "Background" and
% "Related Work" inside this section). If you use subsections at all,
% group them BY TOPIC / APPROACH (one subsection per line of prior work),
% as illustrated below. A short paper can also be a flat section with no
% subsections.
%
% Necessary primer material (terminology, formal definitions, recap of
% prior techniques the reader must know) is woven into the relevant
% topical subsections, NOT separated into its own "Background" subsection.
%
% Reviewers prefer methodological organization (group by *approach*, not
% by *paper*). Avoid the "X et al. did A. Y et al. did B." pattern.

[REPLACE] One short opening paragraph: name the two or three research
communities your work sits between, list the topical subsections that
follow, and point forward to your comparison table if you have one.

\subsection{[Topic / Approach Group 1]}
\label{sec:bg_group1}

\textbf{[Group 1: most closely related approach.]} One line of work
formulates [problem] as [classical formulation]
\cite{PLACEHOLDER_g1_a, PLACEHOLDER_g1_b, PLACEHOLDER_g1_c}. These
methods achieve [property] but assume [assumption]. Weave in any
terminology or formal definitions the reader needs here, rather than
in a separate primer subsection. \method{} follows this line but
differs in [specific difference].

\subsection{[Topic / Approach Group 2]}
\label{sec:bg_group2}

\textbf{[Group 2: alternative approach.]} A second line of work uses
[different formulation, \eg{} ILP, gradient-based, ML-driven]
\cite{PLACEHOLDER_g2_a, PLACEHOLDER_g2_b}. These approaches achieve
[property] at the cost of [limitation]. \method{} differs in
[specific difference, with justification].

\subsection{[Topic / Approach Group 3, if applicable]}
\label{sec:bg_group3}

\textbf{[Group 3: ML / learning-based approach, if applicable.]}
Several recent works apply [deep learning / reinforcement learning / GNN]
to [problem] \cite{PLACEHOLDER_ml_a, PLACEHOLDER_ml_b}.
These methods are orthogonal to \method{} in that [reason for being
orthogonal]; combining them is left as future work.

\subsection{The Gap}
\label{sec:bg_gap}

\textbf{Distinction.} \method{} is the first method, to our knowledge,
that [claim that summarizes the gap in the literature]. The closest prior
method is \cite{PLACEHOLDER_closest}, which differs in [specific
difference]; we compare against it directly in
Section~\ref{sec:experiments}.
